\documentclass{article}

\usepackage{amsmath,amssymb}
\usepackage{xurl}
\usepackage[hidelinks]{hyperref}
\setlength{\emergencystretch}{2em}

\input{command}

\title{Wolf Theorem ``Asymptotic Convergence I'' --- Jordan-Scaling Qualifications}
\author{}
\date{2026-08-23}

\begin{document}
\maketitle
\gapnote{scope-restriction}{open}

\section{The source statement}

Wolf defines $\mu$ as the largest modulus of an eigenvalue of $T$ in the open
unit disc and $d_\mu$ as the largest dimension of a corresponding Jordan
block~\cite[Chapter~8, lines~1221--1272]{Wolf2012Quantum}.  Theorem
``Asymptotic convergence I'' then compares
\begin{equation}
  \lVert\widehat T^n-\widehat T_\varphi^n\rVert_\infty
  \quad\text{with}\quad
  \mu^n n^{d_\mu-1}.
\end{equation}
The local source is
\path{Notes/WolfNoteTexSource/ch08_distance_measures.tex}, lines~1221--1272,
especially Equations~(8.106)--(8.109).

\section{Empty and zero subperipheral spectrum}

If there is no eigenvalue in the open unit disc, neither the printed largest
modulus nor $d_\mu$ is defined.  The formal definition
\leanid{Module.End.subperipheralModulus} totalizes only the first datum by
setting $\mu=0$.  On a nonempty subperipheral spectrum, finiteness proves that
this supremum is an attained maximum strictly below one; see
\leanid{Module.End.exists_norm_eq_subperipheralModulus} and
\leanid{Module.End.subperipheralModulus_lt_one} in
\doclink{QICLean/Analysis/SubperipheralSpectrum.lean}.

The equality $\mu=0$ still does not imply emptiness: it means exactly that
every subperipheral eigenvalue is zero.  A zero-eigenvalue Jordan block of
dimension $D>1$ is nilpotent but its positive powers need not vanish for
$1\le n<D$.  Consequently an upper bound by
$C_2\mu^n n^{d_\mu-1}=0$ cannot hold in that range.  This phenomenon already
occurs for the classical three-state transition which sends state~$1$ to
itself, state~$2$ to state~$1$, and state~$3$ to state~$2$: its
column-stochastic matrix is
\begin{equation}
  P=\begin{pmatrix}1&1&0\\0&0&1\\0&0&0\end{pmatrix}.
\end{equation}
Indeed, $P(e_3-e_1)=e_2-e_1\ne0$ and $P(e_2-e_1)=0$, which exhibits a
size-two Jordan chain at zero.  The associated measure-and-prepare quantum
channel is positive and trace preserving.  Its phase-weighted peripheral map
kills this chain, so at $n=1$ the left side of Equation~(8.107) is nonzero
while its printed right side is zero.  Thus this is not merely a boundary case
outside the theorem's hypotheses.  A qualified zero-modulus statement must
either start at the nilpotence index or treat the finite initial range
separately without the printed factor.  The source's $n\in\mathbb N$ is used
as a positive exponent in this proof; Lean's natural numbers include zero,
whose zeroth power also requires an explicit $0^0$ convention and does not
belong to the argument at source lines~1247--1266.

For $0<\mu\le1$ and $d_\mu-1\le n$, the exact fixed-block consequence of
Equation~(8.104) is formalized as
\leanid{Matrix.l2_opNorm_jordanBlock_pow_polynomial_bounds} in
\doclink{QICLean/Analysis/JordanBlockAsymptotics.lean}.  It proves the
two-sided $\mu^n n^{d_\mu-1}$ scaling for one supplied block; it does not
silently extend that formula across the zero-modulus boundary.

\section{The block-diagonal comparison}

At source line~1261 the notes print
\begin{equation}
  \lVert J^n\rVert_\infty\le \lVert X^n\rVert_\infty,
\end{equation}
where $X$ is one block of the block-diagonal matrix $J$.  In the Hilbert-space
operator norm the inequality for an arbitrary chosen block has the opposite
direction.  The intended eventual-dominance argument must prove that a block
of largest subperipheral modulus and, among ties, largest dimension realizes
the maximum of all block norms.  Only then does one obtain the required
eventual equality $\lVert J^n\rVert_\infty=\lVert X^n\rVert_\infty$.

\section{The Jordan condition infimum}

For a chosen similarity $X=AJA^{-1}$, submultiplicativity gives
\begin{equation}
 (\lVert A\rVert_\infty\lVert A^{-1}\rVert_\infty)^{-1}
 \lVert J^n\rVert_\infty
 \le \lVert X^n\rVert_\infty
 \le
 \lVert A\rVert_\infty\lVert A^{-1}\rVert_\infty
 \lVert J^n\rVert_\infty.
\end{equation}
This exact basis-level comparison is formalized as
\leanid{Matrix.l2_opNorm_pow_similarity_bounds} in
\doclink{QICLean/Analysis/JordanSimilarity.lean}.  Wolf instead writes the
infimum $\kappa_T$ over all Jordan changes of basis directly in Equation
(8.108).  That replacement needs either attainment of the infimum or a
limiting argument compatible with the chosen normalized Jordan form.  No such
result is currently available in the imported Lean Jordan API, so the formal
declaration does not assume a minimizing basis.

\section{Verdict}

\textbf{Verdict: source qualifications and missing Jordan-form
infrastructure.}

The shared modulus, the basis-level similarity inequality, and the
positive-modulus single-block scaling are proved.  The full statement for
$\widehat T^n-\widehat T_\varphi^n$ still requires a constructive Jordan
decomposition, largest-block data $d_\mu$, the direct-sum norm maximum,
eventual dominance, the zero-modulus/positive-exponent qualification, and a
justified passage to $\kappa_T$.  Until those dependencies are supplied,
Equations~(8.106)--(8.109) should not be marked formalized.  Tracking:
\href{https://github.com/LionSR/QICLean/issues/297}{QICLean issue \#297}.

\bibliographystyle{alpha}
\bibliography{references}

\end{document}
